{\color{ink}\bfseries 例 5-6 \(\cdot \) 合订本30 \(\cdot \) 图1}\par \textbf{前缀码树，粉色为已占用叶}\par\vspace{2bp}\begin{tikzpicture}[x=1bp,y=-1bp,>=Stealth,font=\fontsize{8.5}{10.7}\selectfont,every node/.style={inner sep=1bp},line width=.45bp]
\draw[->] (83,15) -- (52,35);
\node[anchor=north west,align=left] at (63.5,17.0){0};
\draw[->] (83,15) -- (115,35);
\node[anchor=north west,align=left] at (95.0,17.0){1};
\draw[->] (52,49) -- (34,71);
\node[anchor=north west,align=left] at (39.0,52.0){0};
\draw[->] (52,49) -- (77,71);
\node[anchor=north west,align=left] at (60.5,52.0){1};
\draw[->] (34,85) -- (24,107);
\node[anchor=north west,align=left] at (25.0,88.0){0};
\draw[->] (34,85) -- (91,107);
\node[anchor=north west,align=left] at (58.5,88.0){1};
\draw[->] (24,121) -- (13,143);
\node[anchor=north west,align=left] at (14.5,124.0){0};
\draw[->] (24,121) -- (46,143);
\node[anchor=north west,align=left] at (31.0,124.0){1};
\draw[->] (91,121) -- (79,143);
\node[anchor=north west,align=left] at (81.0,124.0){0};
\draw[->] (91,121) -- (113,143);
\node[anchor=north west,align=left] at (98.0,124.0){1};
\node[fill=white,draw,anchor=center] at (83,8){根};
\node[fill=white,draw,anchor=center] at (52,42){0};
\node[fill=pink,draw,anchor=center] at (115,42){1};
\node[fill=white,draw,anchor=center] at (34,78){00};
\node[fill=pink,draw,anchor=center] at (77,78){01};
\node[fill=white,draw,anchor=center] at (24,114){000};
\node[fill=white,draw,anchor=center] at (91,114){001};
\node[fill=green,draw,anchor=center] at (13,150){0000};
\node[fill=green,draw,anchor=center] at (46,150){0001};
\node[fill=green,draw,anchor=center] at (79,150){0010};
\node[fill=green,draw,anchor=center] at (113,150){0011};
\end{tikzpicture}\par 剩余00前缀下还有2层，最多容纳\(2^{2}=\allowbreak{}4\)个长度为4的码字。
